\cleardoublepage{}
{
    \chapternonum{毕业设计开题报告、外文翻译的考核}
    \bfseries

    {
        \zihao{4}
        \noindent 导师对开题报告、外文翻译的评语及成绩评定：
    }

    \setlength{\parindent}{2em}

    \par 选题新颖，内容详实，在智能体方面做了很多扎实基础性工作，对拖延症的心理及量化研究方面有一定的前期成果。当前已经取得了较为突出的进展，预期能优秀完成预期成果。

    \designproposaleval[14][4]
    \signature{导师签名}
    % comment the line above and uncomment the line below if you want to set a signature with a specific date.
    % \signaturewithdate{导师签名}{1897}{5}{21}

    {
        \zihao{4}
        \noindent 学院盲审专家对开题报告、外文翻译的评语及成绩评定：
    }

    \setlength{\parindent}{2em}

    \par 该开题报告选题具有较强的现实意义与应用价值，聚焦拖延症干预这一典型认知与行为问题，结合智能体技术提出解决方案，体现出良好的跨学科融合意识。报告结构完整，背景分析充分，对现有任务管理工具与AIGC应用的不足剖析较为深入，理论与实践结合较好。技术路线清晰，系统架构设计较为具体，工程实现路径明确，可行性分析较为充分。但仍存在部分不足：一是文献综述与理论支撑的系统性和批判性有待加强；二是部分技术指标设定较理想化，缺乏实验或数据支撑；三是创新点需进一步凝练与突出。总体达到本科毕业设计开题要求。

    \mbox{} \vfill
    \designproposaleval[13][4]
    \signature{开题报告审核负责人（签名/签章）}
    % comment the line above and uncomment the line below if you want to set a signature with a specific date.
    % \signaturewithdate{开题报告审核负责人（签名/签章）}{1897}{5}{21}
}
